%%% FIELD proof-orthogonal-normal-signal
Fix \(e=(u,c,c')\in\mathcal E_{\mathrm{per}}\).  Introduce the four signed cell indices
\[
\mathcal A_e=\{(c,1),(c,0),(c',1),(c',0)\},\qquad
\sigma_{c,1}=1,\quad \sigma_{c,0}=-1,\quad
\sigma_{c',1}=-1,\quad \sigma_{c',0}=1.
\]
For each \((a,s)\in\mathcal A_e\), the definitions of \(p_{a,s}\) and \(m_{a,s,u}\), iterated conditioning, and positivity of \(\widetilde p_{a,s}\) give
\[
\begin{aligned}
&E_P\!\left[
 \frac{\mathbf 1_{a,s}}{\widetilde p_{a,s}(X)}
 \{\Delta Y_u-\widetilde m_{a,s,u}(X)\}
 \mathrel{\Big|}X\right]\\
&\qquad=
\frac{1}{\widetilde p_{a,s}(X)}
 E_P\!\left[
 \mathbf 1_{a,s}
 E_P\{\Delta Y_u-\widetilde m_{a,s,u}(X)\mid G,S,X\}
 \mathrel{\Big|}X\right]\\
&\qquad=
\frac{p_{a,s}(X)}{\widetilde p_{a,s}(X)}
\{m_{a,s,u}(X)-\widetilde m_{a,s,u}(X)\}.
\end{aligned}
\tag{1}
\]
All divisions in (1) are legitimate because the premise imposes
\(\widetilde p_{a,s}(X)\geq\varepsilon/2>0\) almost surely.

By Definition \(\ref{def:orthogonal-complete-normal-signal}\), the candidate restriction level and its true counterpart are
\[
\widetilde g_e(X)=\sum_{(a,s)\in\mathcal A_e}
 \sigma_{a,s}\widetilde m_{a,s,u}(X),
\qquad
g_e(X)=\sum_{(a,s)\in\mathcal A_e}
 \sigma_{a,s}m_{a,s,u}(X).
\]
Summing (1) with the signs \(\sigma_{a,s}\), adding
\(\widetilde g_e(X)-g_e(X)\), and using the definition of \(S_e\) yields the exact conditional-drift identity
\[
\begin{aligned}
E_P[S_e(O;\widetilde\eta)\mid X]-g_e(X)
&=\sum_{(a,s)\in\mathcal A_e}\sigma_{a,s}
 \{\widetilde m_{a,s,u}(X)-m_{a,s,u}(X)\}
 \left\{1-\frac{p_{a,s}(X)}{\widetilde p_{a,s}(X)}\right\}\\
&=\sum_{(a,s)\in\mathcal A_e}\sigma_{a,s}
 \frac{\{\widetilde p_{a,s}(X)-p_{a,s}(X)\}
       \{\widetilde m_{a,s,u}(X)-m_{a,s,u}(X)\}}
      {\widetilde p_{a,s}(X)}.
\end{aligned}
\tag{2}
\]
Thus the drift is exactly a finite sum of products of a propensity error and an outcome-regression error.  In particular, if only one nuisance is varied from the truth while the other is held at its true value, (2) is identically zero.  More generally, along a Gateaux path
\(\widetilde p_h=p+h\dot p+o(h)\) and
\(\widetilde m_h=m+h\dot m+o(h)\) in the displayed \(L_4(P_X)\) norms, (2) is \(O(h^2)\) in \(L_2(P_X)\).  Its derivative at \(h=0\) is therefore zero, proving Neyman orthogonality with respect to either nuisance at the truth.

The clipping premise gives
\(\widetilde p_{a,s}(X)^{-1}\leq 2/\varepsilon\) almost surely.  Hence the triangle inequality applied to (2), followed by H\"older's inequality with exponents \(4,4,2\), gives
\[
\begin{aligned}
\|E_P[S_e(O;\widetilde\eta)\mid X]-g_e(X)\|_{P_X,2}
&\leq \frac{2}{\varepsilon}
 \sum_{(a,s)\in\mathcal A_e}
 \|\widetilde p_{a,s}-p_{a,s}\|_{P_X,4}
 \|\widetilde m_{a,s,u}-m_{a,s,u}\|_{P_X,4}.
\end{aligned}
\tag{3}
\]
The coordinate set \(\mathcal E_{\mathrm{per}}\), the retained cell set, and the period set are finite.  Therefore, summing the squared coordinate bounds in (3) and using equivalence of norms in the resulting fixed finite-dimensional spaces produces a finite constant \(C_\varepsilon\), depending only on \(\varepsilon\) and those fixed index sets, such that
\[
\|E_P[S^{\mathrm{per}}(O;\widetilde\eta)\mid X]-g(X)\|_{P_X,2}
\leq C_\varepsilon
\|\widetilde p-p\|_{P_X,4}
\|\widetilde m-m\|_{P_X,4}.
\tag{4}
\]
This argument also covers the zero-dimensional convention, in which both sides of (4) are zero.

Finally, substituting the true nuisance \(\eta=(p,m)\) into Definition \(\ref{def:orthogonal-complete-normal-signal}\) gives
\[
S_e(O;\eta)=g_e(X)+Z^{\mathrm{per}}_{e,P}(O).
\tag{5}
\]
Every \(e\in\mathcal E_{\mathrm{per}}\) is, by Definition \(\ref{def:complete-normal-projection}\), a localized observable period-specific equality row.  Under those equalities, Definition \(\ref{def:local-observed-model}\), consistently with Theorem \(\ref{thm:observable-equivalence}\), gives \(g_e(X)=0\) \(P_X\)-almost surely.  Equation (5) consequently yields
\(S_e(O;\eta)=Z^{\mathrm{per}}_{e,P}(O)\) almost surely for every retained row, and hence
\(S^{\mathrm{per}}(O;\eta)=Z_P^{\mathrm{per}}(O)\) almost surely.  This proves all assertions.
